\documentclass[sigconf]{acmart}

%% Remove ACM reference format and copyright for homework submission
\settopmatter{printacmref=false}
\renewcommand\footnotetextcopyrightpermission[1]{}
\pagestyle{plain}

\usepackage{verbatim}
\usepackage{booktabs}

%%
%% Title
%%
\title{Extending Hermes Agent with a Hierarchical Bioinformatics Memory Backend Built on LakeHarbor's Structure-Aware Architecture}

\author{Chan Chi-Ru (詹麒儒)}
\affiliation{%
  \institution{National Taiwan University}
  \country{Taiwan}
}
\email{d12528018@ntu.edu.tw}

%%
%% Abstract
%%
\begin{abstract}
Hermes Agent (NousResearch) is a general-purpose autonomous AI agent centered on a closed learning loop. Its existing memory backend (SQLite FTS5 + Honcho) faces two critical bottlenecks in the bioinformatics domain: computational redundancy caused by the lack of intermediate feature caching, and context explosion from returning full-text blocks to the LLM. Based on the structure-aware data lake architecture of LakeHarbor (ICDE 2024), we propose a three-tier hierarchical memory backend delivered as an MCP Server plugin for Hermes Agent. The system integrates DuckDB for hybrid structured-semantic queries, LLMLingua for semantic text compression in mid-term memory, and DeepSeek-OCR for optical compression in long-term memory archival. The two compression techniques are complementary and non-overlapping, collectively covering both textual and visual bioinformatics outputs without modifying Hermes Agent's core codebase.
\end{abstract}

\begin{document}
\maketitle

%%
%% 1. Problem Definition
%%
\section{Problem Definition}

\textbf{Hermes Agent} employs SQLite FTS5 full-text search as its cross-session memory index and Honcho for user behavior modeling. While effective for general tasks, this design exposes two fundamental deficiencies in bioinformatics scenarios:

\begin{enumerate}
  \item \textbf{Computational Redundancy.} Bioinformatics pipelines (STARsolo alignment, Squidpy spatial clustering) require hours to process a single Visium HD sample ($\sim$50 GB). Hermes's FTS5 memory layer stores only plain-text summaries and cannot cache intermediate feature matrices, causing semantically similar queries to repeatedly trigger expensive heavy computation.

  \item \textbf{Context Explosion \& Token Cost.} Bioinformatics analysis reports can contain hundreds of thousands of words. When Hermes's skill system returns full-text blocks, the cost using GPT-4o (\$5/1M input tokens) is prohibitive and frequently exceeds the LLM context window.
\end{enumerate}

Existing data lake systems such as LakeHarbor, while elevating structured schemas to first-class citizens, still serve all data at a \textbf{static, homogeneous resolution}. They do not address the core Agent-driven question of \textit{at what granularity should memory be retained}. Furthermore, Hermes Agent currently has no plugin integrating any data lakehouse architecture, representing a clear engineering gap.

%%
%% 2. Prior Work & Research Gap
%%
\section{Prior Work \& Research Gap}

\textbf{Primary Reference:} \textit{LakeHarbor: Making Structures First-Class Citizens in Data Lakes}~\cite{lakeharbor2024} (ICDE 2024, pp.\ 5583--5592).

LakeHarbor proposes elevating structured schemas to first-class citizens within data lakes, enabling lakes to retain schema-on-read flexibility while achieving the query performance of structured warehouses. Its core contribution is a unified structure-aware storage layer that bridges the gap between data lakes and data warehouses.

\textbf{Research Gap.} LakeHarbor assumes all data is presented to every querier at the same resolution. In an LLM Agent-driven bioinformatics scenario, this creates two unresolved problems:

\begin{itemize}
  \item The system cannot automatically \textbf{downgrade-compress} data based on temporal distance or query importance (full text $\rightarrow$ summary $\rightarrow$ index).
  \item There is no mechanism to embed context compression into the storage engine layer, leaving all token cost management entirely to the application layer.
\end{itemize}

Additionally, agent-first data system design~\cite{aioverlords2025} argues that future databases must natively support agent workloads, yet no concrete implementation exists for bioinformatics vertical domains.

%%
%% 3. Proposed Solution
%%
\section{Proposed Solution}

We propose delivering a bioinformatics-specific memory backend as an \textbf{MCP Server plugin} for Hermes Agent, extending LakeHarbor's structure-aware storage design with a \textbf{Multi-Resolution Semantic Cache}. Hermes's \texttt{mcp\_serve.py} mechanism enables seamless plugin integration; the Agent accesses the new backend via MCP Tool Calls without any modification to core code.

\subsection{Three-Tier Medallion Architecture}

The system adopts a three-tier Medallion architecture with lazy execution and intelligent fallback:

\begin{verbatim}
Hermes Agent Query
    | (MCP Tool Call: bio_memory_query)
    v
[L1 Gold]  Multi-Res Cache --hit--> Return (cost~0)
    | miss
    v
[L2 Silver] DuckDB Store --hit--> LLMLingua -> Append L1
    | miss
    v
[L3 Bronze] Raw Lake -> Pipeline -> Write L2 -> L1
\end{verbatim}

\subsection{Layer Descriptions}

\textbf{Bronze Layer (L3) --- Immutable Raw Lake.}
Stores raw FASTQ files and image matrices. Immutable by design. Heavy pipelines (STARsolo, Squidpy) are triggered only when L2 lacks the required features, with results written back through L2 and L1.

\textbf{Silver Layer (L2) --- DuckDB Structured Feature Store.}
Extends LakeHarbor's structure-aware design by storing schema-aware \texttt{.h5ad}-extracted count matrices in Parquet format. The DuckDB \texttt{vss} extension enables hybrid ``gene name structured filtering + semantic vector search'' queries, directly addressing the limitations of Hermes's existing SQLite FTS5 backend~\cite{duckdb2019}.

\textbf{Gold Layer (L1) --- Multi-Resolution Semantic Cache.}
This is the core innovation beyond LakeHarbor's static storage model. It automatically manages Hermes's bioinformatics memory across three resolutions based on TTL and access frequency, replacing Honcho + FTS5 as the domain-specific memory layer:

\begin{enumerate}
  \item \textit{Recent Memory (high-res, TTL = 7 days):} Complete plain-text analysis reports for full Agent reasoning.

  \item \textit{Mid-term Memory (mid-res, TTL = 90 days):} \textbf{LLMLingua semantic compression}~\cite{llmlingua2023} at 20$\times$ ratio, removing redundant words while preserving complete semantic structure (DEG lists, statistical conclusions) for Agent reasoning without requiring full text.

  \item \textit{Long-term Memory (low-res, permanent):} \textbf{DeepSeek-OCR optical compression}~\cite{deepseekocr2025}---encodes aged reports (including text and figures) into minimal vision tokens (100 tokens reconstructing 800+ text tokens at 96.8\% precision), using the visual modality as a compression intermediary. Used for indexing only, not reasoning. Full reports are expanded only on semantic hit (cosine similarity $\geq$ 0.88). Future work will explore whether DeepSeek-OCR can fully replace LLMLingua.
\end{enumerate}

The two compression techniques are fundamentally complementary: LLMLingua operates at the \textit{semantic layer} (preserving structure for reasoning), while DeepSeek-OCR operates at the \textit{visual layer} (maximizing compression for indexing).

%%
%% 4. Experiments & Verification
%%
\section{Experiments \& Verification}

We implement a prototype using the 10x Genomics public PBMC Visium HD dataset ($\sim$45 GB), building an MCP Server with DuckDB (\texttt{vss} extension) and LLMLingua middleware, connected to a real Hermes Agent instance.

\textbf{Experimental Setup.} Simulate Hermes Agent performing 100 consecutive exploratory queries targeting the tumor microenvironment.

\textbf{Evaluation Metrics:}
\begin{enumerate}
  \item \textit{Token Cost Analysis:} Compare total API token consumption between Hermes native FTS5 memory vs.\ this module's L1 cache hit.
  \item \textit{Latency Benchmark:} Measure end-to-end latency for L1 hit ($<$1\,s) vs.\ L2 retrieval ($\sim$30\,s) vs.\ L3 recomputation ($\sim$4\,h), verified with paired $t$-test ($\alpha = 0.05$).
  \item \textit{Semantic Retention Rate:} Use BERTScore F1 to measure consistency between compressed conclusions and full-data results, targeting $\geq 0.92$.
\end{enumerate}

%%
%% 5. Expected Results
%%
\section{Expected Results}

Based on LLMLingua's 20$\times$ compression ratio and Prefix KV Cache's $\sim$90\% token savings, we expect this extension module to achieve \textbf{$\geq$80\% API token reduction} on redundant bioinformatics queries, and prevent over 70\% of subsequent queries from triggering L3 heavy recomputation. This work will demonstrate that LakeHarbor's structure-aware architecture can serve as the engineering foundation for general AI Agent memory backends, and that embedding context compression at the storage layer is the key path to scalable, agent-driven bioinformatics systems.

%%
%% References
%%
\bibliographystyle{ACM-Reference-Format}
\begin{thebibliography}{9}

\bibitem{lakeharbor2024}
Hiroyuki Yamada, Masaru Kitsuregawa, and Kazuo Goda. 2024.
\newblock LakeHarbor: Making Structures First-Class Citizens in Data Lakes.
\newblock In \textit{Proc.\ 40th IEEE International Conference on Data Engineering (ICDE '24)}, pp.\ 5583--5592.

\bibitem{hermesagent2025}
NousResearch. 2025.
\newblock \textit{Hermes Agent: The agent that grows with you}.
\newblock GitHub. \url{https://github.com/nousresearch/hermes-agent}

\bibitem{aioverlords2025}
Shu Liu et al. 2025.
\newblock Supporting Our AI Overlords: Redesigning Data Systems to be Agent-First.
\newblock \textit{arXiv preprint arXiv:2509.00997}.

\bibitem{llmlingua2023}
Huiqiang Jiang, Qianhui Wu, Chin-Yew Lin, Yuqing Yang, and Lili Qiu. 2023.
\newblock LLMLingua: Compressing Prompts for Accelerated Inference of Large Language Models.
\newblock In \textit{Proc.\ EMNLP 2023}, pp.\ 13358--13376.

\bibitem{deepseekocr2025}
Haoran Wei, Yaofeng Sun, and Yukun Li. 2025.
\newblock DeepSeek-OCR: Contexts Optical Compression.
\newblock \textit{arXiv preprint arXiv:2510.18234}.
\newblock \url{https://github.com/deepseek-ai/DeepSeek-OCR}

\bibitem{duckdb2019}
Mark Raasveldt and Hannes M{\"u}hleisen. 2019.
\newblock DuckDB: An Embeddable Analytical Database.
\newblock In \textit{Proc.\ SIGMOD 2019}, pp.\ 1981--1984.

\end{thebibliography}

\end{document}
